\section*{Limitations}

\paragraph{Methodological scope.}
MH-DIF detects only uniform DIF; logistic regression reveals an additional 24.8\% non-uniform DIF, making 31\% conservative. More fundamentally, DIF detects \emph{that} items function differently but cannot determine \emph{why}---the signal may reflect contamination, capability change, or training methodology shifts. A definitive causal decomposition remains out of reach without access to model training data.

\paragraph{Matching boundary and sample size.}
Cross-subject matching assumes contamination affects a minority of domains; when contamination spans ${>}10$ subjects (${>}18\%$ of MMLU), $\Delta\mathrm{FPR}$ rises to 0.026--0.069. The GSM1K correlation analysis ($N = 33$) has limited statistical power and is exploratory; replication with larger model samples is warranted.

\paragraph{Generalizability.}
All six benchmarks use multiple-choice or binary scoring; generalizability to open-ended generation (e.g., AlpacaEval), coding (e.g., HumanEval), or partial-credit benchmarks is untested and would require polytomous DIF methods. Temporal cohorts are defined by calendar year, a coarse grouping; finer-grained definitions (e.g., by training data cutoff or alignment method) might reveal more targeted patterns.
